\section{Base Bridge Count}

\subsection{Material carrier contribution}

The material side of the electromagnetic bridge is the compact fiber
\[
F=Z_4\times Z_4\times Z_8.
\]
Therefore the material carrier contributes
\[
|F|=4\cdot4\cdot8=128.
\]
This count is fixed once the compact fiber is fixed.

\subsection{Radiative/vibrational bridge-placement contribution}

The radiative/vibrational bridge-placement carrier contributes
\[
|V_{\rm rad}^{\rm bridge}|=3^2=9.
\]
This count is not a magnetic-sector count and is not inserted to force
the number \(137\). Its provenance is radiative/vibrational bridge
placement.

A one-dimensional scalar-displacement relation with a distinguished datum
has the qualitative sign-status quotient
\[
\{-,0,+\},
\]
corresponding to position below, at, or above the datum. The displacement
algebra itself is not reduced to three elements; rather, the
pre-projection placement readout has three sign-status classes.

Radiation is treated here as a one-dimensional vibratory seed. In
three-dimensional placement geometry, the orientation space of a
one-dimensional axis is two-dimensional:
\[
\dim \operatorname{Gr}(1,3)=2.
\]
Thus the electromagnetic bridge reads the radiative/vibrational placement
over two bridge-facing placement dimensions. The radiation-sector paper
treats this as the vibrational placement layer of the radiation
architecture \cite{Wells2026Radiation}.

Given the two-dimensional placement readout and the ternary sign-status
quotient in each placement dimension,
\[
V_{\rm rad}^{\rm bridge}
\cong
\{-,0,+\}^2,
\]
and therefore
\[
|V_{\rm rad}^{\rm bridge}|
=
|\{-,0,+\}^2|
=
3^2
=
9.
\]
Only the placement count is used in the alpha derivation. No group law on
\(V_{\rm rad}^{\rm bridge}\) is required.

\subsection{Disjoint-union access count}

The base bridge count uses the carrier-access composition rule:
\[
\boxed{
\text{mechanism-distinct, independently sufficient access routes add;}
}
\]
\[
\boxed{
\text{simultaneously required labels multiply;}
}
\]
\[
\boxed{
\text{mere refinements or marginals of the same mechanism collapse rather than add.}
}
\]

This rule separates three cases. If two labels must be specified
simultaneously to identify a single elementary event, the correct access
space is a Cartesian product. If two structures supply different-kind,
independently sufficient access routes to the same scalar bridge
observable, the correct access space is a disjoint union. If one
structure is only a refinement or marginalization of the other, it is not
counted as an additional independent access route.

The electromagnetic bridge is a carrier-level material/radiative
interface. It is not a joint state space whose elements are ordered pairs
\[
(f,v)\in F\times V_{\rm rad}^{\rm bridge}.
\]
A product count would give
\[
|F\times V_{\rm rad}^{\rm bridge}|
=
|F|\,|V_{\rm rad}^{\rm bridge}|.
\]
That would be the correct count only if an electromagnetic bridge event
required simultaneous specification of one material element and one
radiative-placement element as a joint state label.

The bridge coupling derived here is instead a scalar carrier-level
normalization. Material carrier access and radiative/vibrational
bridge-placement access are mechanism-distinct and independently
sufficient at the bridge level. They are not mutually indexed degrees of
freedom of a single joint state, and neither is merely a marginal
refinement of the other.

Thus the relevant access space is the disjoint union
\[
F\sqcup V_{\rm rad}^{\rm bridge},
\]
not the tensor/product-state access space
\[
F\times V_{\rm rad}^{\rm bridge}.
\]
The base bridge count is therefore
\[
B
=
|F\sqcup V_{\rm rad}^{\rm bridge}|
=
|F|+|V_{\rm rad}^{\rm bridge}|.
\]
Substituting the two carrier counts gives
\[
B=128+9=137.
\]

\subsection{Base count result}

The base electromagnetic bridge count is
\[
\boxed{
B=137.
}
\]
This is the zeroth-order compact-fiber bridge capacity before
cover-mediated feedback is included. It is not the final value of
\[
x=\alpha_F^{-1}.
\]
The final effective bridge capacity is obtained by adding feedback
self-consistently,
\[
x=B+\frac{K}{x},
\]
and then including the closed-return correction
\[
-\frac{K N_m}{(N_e-1)x^3}.
\]

The dependency structure is:
\[
F=Z_4\times Z_4\times Z_8
\quad\Rightarrow\quad
|F|=128,
\]
\[
\{-,0,+\}\ \text{over two radiative placement dimensions}
\quad\Rightarrow\quad
|V_{\rm rad}^{\rm bridge}|=3^2=9,
\]
\[
\text{mechanism-distinct disjoint-union access}
\quad\Rightarrow\quad
B=|F|+|V_{\rm rad}^{\rm bridge}|.
\]
Thus
\[
\boxed{
B=137
}
\]
is a structural base count determined by carrier access and
bridge-placement readout, not a fitted numerical input.